% ID: SEL-D0002-OB-E01
% ESTADO: validada. S003, 2026-11-07; tres pasadas del mismo autor, comprobación separada y revisión didáctica/visual.
\textbf{a) Variables, restricciones y región.}\par\nopagebreak[4]
Sean $x$ e $y$ los kg de concentrado A y B. Son cantidades continuas y no negativas; no se exige integridad. Cada disponibilidad es una cota superior:
\[
\begin{cases}
4.5x+7.5y\leq67.5&\text{(Colombia)},\\
3x\leq30&\text{(Etiopía)},\\
1.5y\leq9&\text{(Costa Rica)},\\
x\geq y/2,\quad x\geq0,\quad y\geq0.&
\end{cases}
\quad\Longleftrightarrow\quad
\begin{cases}
3x+5y\leq45,\\x\leq10,\\y\leq6,\\y\leq2x,\\x\geq0,\ y\geq0.
\end{cases}
\]
La región es la intersección de esos semiplanos, incluidas sus fronteras. Sus vértices en orden son
\[
\boxed{(0,0),\ (10,0),\ (10,3),\ (5,6),\ (3,6).}
\]
Se obtienen intersectando las fronteras consecutivas: $y=0$ con $y=2x$; $x=10$ con $y=0$; $x=10$ con $3x+5y=45$; $y=6$ con $3x+5y=45$; y $y=6$ con $y=2x$. Todos cumplen las restricciones; las demás intersecciones no factibles se descartan.
\begin{center}\includegraphics[width=0.67\linewidth]{figuras/s003-region-factible.pdf}\end{center}

\textbf{b) Comprobación de $(7,5)$.}\par\nopagebreak[4]
No es posible: necesitaría $4.5(7)+7.5(5)=69$ kg de Colombia, más de los 67.5 disponibles. Las otras restricciones sí se cumplen ($21\leq30$, $7.5\leq9$ y $7\geq2.5$), pero incumplir una sola ya excluye el punto.

\textbf{c) Beneficio máximo.}\par\nopagebreak[4]
La función objetivo es $G=2x+4y$ euros. Una función lineal sobre este polígono cerrado y acotado alcanza máximo en un vértice:
\[
\begin{array}{c|ccccc}
(x,y)&(0,0)&(10,0)&(10,3)&(5,6)&(3,6)\\\hline
G&0&20&32&34&30
\end{array}
\]
Por tanto, producir \boxed{x=5\text{ kg de A},\quad y=6\text{ kg de B}} proporciona \boxed{34\text{ euros}}. Consume 67.5 kg de Colombia, 15 de Etiopía y 9 de Costa Rica. Comprobación global alternativa:
\[
G=\frac23(3x+5y)+\frac23y\leq\frac23(45)+\frac23(6)=34.
\]
La igualdad exige simultáneamente $3x+5y=45$ e $y=6$, que determinan únicamente $(5,6)$.
